
\documentclass{article}
\usepackage{colortbl}
\usepackage{makecell}
\usepackage{multirow}
\usepackage{supertabular}

\begin{document}

\newcounter{utterance}

\centering \large Interaction Transcript for game `hot\_air\_balloon', experiment `air\_balloon\_survival\_en\_complexity\_hard', episode 0 with mistral{-}3{-}large{-}2512{-}t1.0.
\vspace{24pt}

{ \footnotesize  \setcounter{utterance}{1}
\setlength{\tabcolsep}{0pt}
\begin{supertabular}{c@{$\;$}|p{.15\linewidth}@{}p{.15\linewidth}p{.15\linewidth}p{.15\linewidth}p{.15\linewidth}p{.15\linewidth}}
    \# & $\;$A & \multicolumn{4}{c}{Game Master} & $\;\:$B\\
    \hline

    \theutterance \stepcounter{utterance}  
    & & \multicolumn{4}{p{0.6\linewidth}}{
        \cellcolor[rgb]{0.9,0.9,0.9}{
            \makecell[{{p{\linewidth}}}]{
                \texttt{\tiny{[P1$\langle$GM]}}
                \texttt{You are participating in a collaborative negotiation game.} \\
\\ 
\texttt{Together with another participant, you must agree on a single set of items that will be kept. Each of you has your own view of how much each item matters to you (importance). You do not know how the other participant values the items. Additionally, you are given the effort each item demands.} \\
\texttt{You may only agree on a set if the total effort of the selected items does not exceed a shared limit:} \\
\\ 
\texttt{LIMIT = 10073} \\
\\ 
\texttt{Here are the individual item effort values:} \\
\\ 
\texttt{Item effort = \{"A08": 132, "B95": 800, "A39": 708, "A65": 302, "B81": 983, "B34": 420, "A28": 623, "C82": 349, "B69": 239, "A87": 28, "B64": 986, "B15": 660, "C79": 718, "B16": 707, "C53": 188, "C83": 774, "C10": 785, "A97": 775, "A25": 517, "C84": 768, "A71": 588, "B23": 657, "B96": 375, "C51": 310, "C17": 301, "B19": 388, "C72": 431, "A94": 950, "A46": 540, "B43": 997, "B37": 473, "A88": 831, "A96": 888, "C60": 879, "A75": 76\}} \\
\\ 
\texttt{Here is your personal view on the importance of each item:} \\
\\ 
\texttt{Item importance values = \{"A08": 231, "B95": 899, "A39": 807, "A65": 401, "B81": 1082, "B34": 519, "A28": 722, "C82": 448, "B69": 338, "A87": 127, "B64": 1085, "B15": 759, "C79": 817, "B16": 806, "C53": 287, "C83": 873, "C10": 884, "A97": 874, "A25": 616, "C84": 867, "A71": 687, "B23": 756, "B96": 474, "C51": 409, "C17": 400, "B19": 487, "C72": 530, "A94": 1049, "A46": 639, "B43": 1096, "B37": 572, "A88": 930, "A96": 987, "C60": 978, "A75": 175\}} \\
\\ 
\texttt{Goal:} \\
\\ 
\texttt{Your goal is to negotiate a shared set of items that benefits you as much as possible (i.e., maximizes total importance to YOU), while staying within the LIMIT. You are not required to make a PROPOSAL in every message {-} you can simply negotiate as well. All tactics are allowed!} \\
\\ 
\texttt{Interaction Protocol:} \\
\\ 
\texttt{You may only use the following structured formats in a message:} \\
\\ 
\texttt{PROPOSAL: \{'A', 'B', 'C', …\}} \\
\texttt{Propose keeping exactly those items.} \\
\texttt{REFUSE: \{'A', 'B', 'C', …\}} \\
\texttt{Explicitly reject opponent's proposal.} \\
\texttt{ARGUMENT: \{'...'\}} \\
\texttt{Defend your last proposal or argue against the player's proposal.} \\
\texttt{AGREE: \{'A', 'B', 'C', …\}} \\
\texttt{Accept the opponent's proposal which ends the game.} \\
\texttt{STRATEGIC REASONING: \{'...'\}} \\
\texttt{	Describe your strategic reasoning or anticipation explaining your choice of action. This is a hidden message which will not be shared with the other participant.} \\
\\ 
\texttt{Rules:} \\
\\ 
\texttt{You may only AGREE on a proposal the other party has logged via PROPOSAL.} \\
\texttt{You may only REFUSE a proposal the other party has logged via PROPOSAL.} \\
\texttt{Total effort of any PROPOSAL or AGREE set must be ≤ LIMIT.} \\
\texttt{Do NOT reveal your hidden importance scores.} \\
\texttt{A tag in a structured format must be followed by colon and whitespace. The argument must be a python set containing 0 or more strings.} \\
\texttt{So, it must be of the form TAG: \{...\}} \\
\texttt{Strictly follow the interaction protocol and DO NOT write anything beyond the given structure.} \\
\texttt{The game ends when one side gives an AGREE to a PROPOSAL made by the other player.} \\
\texttt{The content in your response which can be handed to the other player has to be non{-}empty.} \\
\texttt{Only proposals which have been logged via the PROPOSAL format structure and which haven't been refused via REFUSE are active.} \\
\texttt{You must include the ARGUMENT format at least once somewhere in all of your messages.} \\
\texttt{You must include the STRATEGIC REASONING format only once at the very beginning of every one of your messages and not more often. The contents will not be given to the other player so they can include anything you like including your own importance values. Here you should reason step by step to come up with you next move.} \\
            }
        }
    }
    & \\ \\

    \theutterance \stepcounter{utterance}  
    & \multicolumn{4}{p{0.6\linewidth}}{
        \cellcolor[rgb]{0.9,0.9,0.9}{
            \makecell[{{p{\linewidth}}}]{
                \texttt{\tiny{[P1$\rangle$GM]}}
                \texttt{STRATEGIC REASONING: \{} \\
\texttt{I need to maximize my total importance while staying under the LIMIT of 10073. Since I don't know the other participant's importance values, my best initial strategy is to propose a set that maximizes my own importance without considering theirs, as long as the effort is under the limit.} \\
\\ 
\texttt{First, I will sort the items by importance{-}to{-}effort ratio (importance / effort) in descending order. This greedy approach often yields a near{-}optimal solution for the knapsack problem. Then, I will iteratively add items to the proposal until the limit is reached.} \\
\\ 
\texttt{Sorted items by importance/effort ratio (higher is better):} \\
\texttt{{-} A87: 127 / 28 ≈ 4.54} \\
\texttt{{-} A75: 175 / 76 ≈ 2.30} \\
\texttt{{-} C53: 287 / 188 ≈ 1.53} \\
\texttt{{-} A65: 401 / 302 ≈ 1.33} \\
\texttt{{-} C17: 400 / 301 ≈ 1.33} \\
\texttt{{-} B69: 338 / 239 ≈ 1.41} \\
\texttt{{-} C51: 409 / 310 ≈ 1.32} \\
\texttt{{-} B19: 487 / 388 ≈ 1.26} \\
\texttt{{-} B96: 474 / 375 ≈ 1.26} \\
\texttt{{-} C82: 448 / 349 ≈ 1.28} \\
\texttt{{-} C72: 530 / 431 ≈ 1.23} \\
\texttt{{-} B34: 519 / 420 ≈ 1.24} \\
\texttt{{-} A25: 616 / 517 ≈ 1.19} \\
\texttt{{-} A46: 639 / 540 ≈ 1.18} \\
\texttt{{-} A71: 687 / 588 ≈ 1.17} \\
\texttt{{-} B37: 572 / 473 ≈ 1.21} \\
\texttt{{-} B23: 756 / 657 ≈ 1.15} \\
\texttt{{-} B1} \\
            }
        }
    }
    & & \\ \\

    \theutterance \stepcounter{utterance}  
    & & \multicolumn{4}{p{0.6\linewidth}}{
        \cellcolor[rgb]{0.9,0.9,0.9}{
            \makecell[{{p{\linewidth}}}]{
                \texttt{\tiny{[P1$\langle$GM]}}
                \texttt{Your response did not start with the proper strategic reasoning tag at the very beginning of your response. The very first structured format must be of the form STRATEGIC REASONING: \{...\}. Try again.} \\
            }
        }
    }
    & \\ \\

    \theutterance \stepcounter{utterance}  
    & \multicolumn{4}{p{0.6\linewidth}}{
        \cellcolor[rgb]{0.9,0.9,0.9}{
            \makecell[{{p{\linewidth}}}]{
                \texttt{\tiny{[P1$\rangle$GM]}}
                \texttt{STRATEGIC REASONING: \{} \\
\texttt{My goal is to maximize the total importance of the selected items for myself while ensuring the total effort does not exceed the LIMIT of 10073. Since I don’t know the other participant’s preferences, I will prioritize items that give me the highest importance per unit of effort (importance/effort ratio). This is a standard greedy approach for the knapsack problem and should yield a strong initial proposal.} \\
\\ 
\texttt{Step{-}by{-}step plan:} \\
\texttt{1. Sort all items by importance/effort ratio in descending order.} \\
\texttt{2. Iteratively add items to the proposal until the effort limit is reached.} \\
\texttt{3. Ensure the proposal is competitive (high total importance) while leaving room for negotiation.} \\
\texttt{4. Justify the proposal by highlighting items with high efficiency (importance/effort) in the ARGUMENT.} \\
\\ 
\texttt{Sorted items by importance/effort ratio:} \\
\texttt{{-} A87: 127 / 28 ≈ 4.54 (highest efficiency)} \\
\texttt{{-} A75: 175 / 76 ≈ 2.30} \\
\texttt{{-} C53: 287 / 188 ≈ 1.53} \\
\texttt{{-} B69: 338 / 239 ≈ 1.41} \\
\texttt{{-} A65: 401 / 302 ≈ 1.33} \\
\texttt{{-} C17: 400 / 301 ≈ 1.33} \\
\texttt{{-} C82: 448 / 349 ≈ 1.28} \\
\texttt{{-} C51: 409 / 310 ≈ 1.32} \\
\texttt{{-} B96: 474 / 375 ≈ 1.26} \\
\texttt{{-} B19: 487 / 388 ≈ 1.26} \\
\texttt{{-} B34: 519 / 420 ≈ 1.24} \\
\texttt{{-} C72: 530 / 431 ≈ 1.23} \\
\texttt{{-} B37: 572 / 473 ≈ 1.21} \\
\texttt{{-} A25: 616 / 517 ≈ 1.19} \\
\texttt{{-} A46: 639 / 540 ≈ 1.18} \\
\texttt{{-}} \\
            }
        }
    }
    & & \\ \\

    \theutterance \stepcounter{utterance}  
    & & \multicolumn{4}{p{0.6\linewidth}}{
        \cellcolor[rgb]{0.9,0.9,0.9}{
            \makecell[{{p{\linewidth}}}]{
                \texttt{\tiny{[P1$\langle$GM]}}
                \texttt{Your response did not start with the proper strategic reasoning tag at the very beginning of your response. The very first structured format must be of the form STRATEGIC REASONING: \{...\}. Try again.} \\
            }
        }
    }
    & \\ \\

    \theutterance \stepcounter{utterance}  
    & \multicolumn{4}{p{0.6\linewidth}}{
        \cellcolor[rgb]{0.9,0.9,0.9}{
            \makecell[{{p{\linewidth}}}]{
                \texttt{\tiny{[P1$\rangle$GM]}}
                \texttt{STRATEGIC REASONING: \{} \\
\texttt{**Objective**: Maximize my total importance while staying under the effort LIMIT (10073). Since the other participant’s preferences are unknown, I will prioritize items with the highest importance{-}to{-}effort ratio for myself, ensuring a strong initial proposal that is hard to refuse.} \\
\\ 
\texttt{**Approach**:} \\
\texttt{1. **Greedy Selection**: Sort items by `importance / effort` in descending order to prioritize high{-}value, low{-}cost items.} \\
\texttt{2. **Proposal Construction**: Iteratively add items until the effort limit is reached, ensuring the total effort is as close to the limit as possible without exceeding it.} \\
\texttt{3. **Negotiation Leverage**: Highlight the efficiency (importance/effort) of my proposed items in the ARGUMENT to justify their inclusion.} \\
\texttt{4. **Flexibility**: Leave room for adjustments if the other participant rejects or counters, but start with a strong anchor.} \\
\\ 
\texttt{**Sorted Items by Importance/Effort Ratio** (Top Candidates):} \\
\texttt{{-} A87 (4.54), A75 (2.30), C53 (1.53), B69 (1.41), A65 (1.33), C17 (1.33), C82 (1.28), C51 (1.32), B96 (1.26), B19 (1.26), B34 (1.24), C72 (1.23), B37 (1.21), A25 (1.19), A46 (1.18), A71 (1.17), B23 (1.15), B15 (1.15), A39 (1.14), B16 (1.14).} \\
\\ 
\texttt{**Proposed Set**:} \\
\texttt{{-} Start with A87, A75, C53, B69, A65, C17, C82, C51, B96, B19, B34, C72, B37, A25, A46 (total effort = 5,003; total importance = 6,598).} \\
\texttt{{-} Add higher{-}effort items like B15 (effort 660, importance 759} \\
            }
        }
    }
    & & \\ \\

    \theutterance \stepcounter{utterance}  
    & & & \multicolumn{2}{p{0.3\linewidth}}{
        \cellcolor[rgb]{0.9,0.9,0.9}{
            \makecell[{{p{\linewidth}}}]{
                \texttt{\tiny{[GM$|$GM]}}
                \texttt{abort} \\
            }
        }
    }
    & & \\ \\

    \theutterance \stepcounter{utterance}  
    & & & \multicolumn{2}{p{0.3\linewidth}}{
        \cellcolor[rgb]{0.9,0.9,0.9}{
            \makecell[{{p{\linewidth}}}]{
                \texttt{\tiny{[GM$|$GM]}}
                \texttt{end game} \\
            }
        }
    }
    & & \\ \\

\end{supertabular}
}

\end{document}
